% Part: turing-machines
% Chapter: undecidability
% Section: universal-tm

\documentclass[../../../include/open-logic-section]{subfiles}

\begin{document}

% SEGMENT OLP-0267-S01

\olfileid{tur}{und}{uni}
\olsection{அனைத்தளாவிய டியூரிங் பொறிகள்}

\olref[enu]{sec} இல் ஒவ்வொரு டியூரிங் பொறியையும் முழுக்களின் ஒரு
முடிவுறு தொடரால் எவ்வாறு விவரிக்கலாம் என்று விவாதித்தோம். இத்தொடர்
டியூரிங் பொறியின் நிலைகள், குறியெழுத்துத் தொகுதி, தொடக்க நிலை,
கட்டளைகள் ஆகியவற்றைக் குறியீடாக்குகிறது. இவ்விவரிப்புகள் அனைத்தின்
கணமும் !!{enumerable} என்றும் சுட்டிக்காட்டினோம். அத்தகைய
விவரிப்புகளின் கணம் !!{denumerable} என்பதால்,~$\Nat$ இலிருந்து
இவ்விவரிப்புகளுக்குச் செல்லும் !!a{surjective} சார்பு உள்ளது என்பது
இதன் பொருள். எடுத்துக்காட்டாக, கான்டரின் வளைந்தோடும் முறையைப்
பயன்படுத்தி அத்தகைய !!a{surjective} சார்பைப் பெறலாம். டியூரிங்
பொறிகளின் எல்லா விவரிப்புகளையும் பட்டியலிடுவதற்கான ஒரு வழியை அது
தருகிறது. அத்தகைய ஓர் எண்ணிடலைத் தேர்ந்துவைத்தால், $1$-ஆம், $2$-ஆம்,
\dots, $e$-ஆம் டியூரிங் பொறியைப் பற்றிப் பேசுவதற்கு இப்போது பொருள்
உண்டு. இவ்வெண்கள் \emph{சுட்டெண்கள்} எனப்படுகின்றன.

\begin{defn}
நாம் தேர்ந்துவைத்த எண்ணிடலில் $M$ ஆனது $e$-ஆம் டியூரிங் பொறி எனில்,
~$e$ ஐ~$M$ இன் ஒரு \emph{சுட்டெண்} என்று கூறுகிறோம். $e$-ஆம்
டியூரிங் பொறிக்கு $M_e$ என்று எழுதுகிறோம்.
\end{defn}

ஒரு பொறிக்கு ஒன்றுக்கு மேற்பட்ட சுட்டெண்கள் இருக்கலாம். எடுத்துக்காட்டாக,
~$M$ இன் இரு விவரிப்புகள் அதன் கட்டளைகளை நாம் பட்டியலிடும் வரிசையில்
வேறுபடலாம்; இவ்வேறு விவரிப்புகளுக்கு வெவ்வேறு சுட்டெண்கள் இருக்கும்.

சுட்டெண்ணிலிருந்து~$M$ இன் விவரிப்பை விளைவுறக் கணிக்கவும், ஒரு
பொறி~$M$ இன் விவரிப்பிலிருந்து அதன் ஒரு சுட்டெண்ணை விளைவுறக் கணிக்கவும்
முடியும் வகையில் டியூரிங் பொறி விவரிப்புகளின் எண்ணிடலைத் தருவது
சாத்தியம் என்பது முக்கியமானது. சர்ச்--டியூரிங் முன்வைப்பின்படி,
சுட்டெண்~$e$ கொண்ட டியூரிங் பொறியின் விவரிப்பை மீட்டெடுத்து, அதற்குரிய
விவரிப்பை வெளியீடாகத் தனது நாடாவில் எழுதும் ஒரு டியூரிங் பொறியைக்
கண்டறிவது அப்போது சாத்தியம். அவ்விவரிப்பு~$\TMstroke$-களின்
தொகுதிகளைக் கொண்ட ஒரு தொடராக இருக்கும் (~$M_e$ ஐ விவரிக்கும் தொடரில்
உள்ள நேர்ம முழுக்களைக் குறிக்கும்).

இது கொடுக்கப்பட்டதும், டியூரிங் பொறிச் சுட்டெண்களின் எந்தச் சார்புகள்
டியூரிங் பொறிகளால் கணிக்கத்தக்கவை? டியூரிங் பொறிச் சுட்டெண்களின்
எந்தப் பண்புகளை டியூரிங் பொறிகளால் தீர்மானிக்கலாம்? என்று கேட்பது
இப்போது இயல்பாகிறது. ஓர் எடுத்துக்காட்டு: சுட்டெண்~$e$ ஐ, சுட்டெண்~$e$
கொண்ட டியூரிங் பொறியின் நிலைகளின் எண்ணிக்கைக்கு அனுப்பும் சார்பு ஒரு
டியூரிங் பொறியால் கணிக்கத்தக்கது. அத்தகைய டியூரிங் பொறி செய்வது இதுதான்:
$e$ எண்ணிக்கையிலான~$\TMstroke$-களின் ஒரே தொகுதியைக் கொண்ட நாடாவில்
தொடங்கப்பட்டால், முதலில் $e$ ஐ அதன் விவரிப்பாகக் குறிவிலக்கும்.
அவ்விவரிப்பு இப்போது நாடாவில்~$\TMstroke$-களின் தொகுதிகளைக் கொண்ட
ஒரு தொடரால் குறிப்பிடப்படுகிறது. இத்தொடரின் முதல் !!{element} ஆனது
நிலைகளின் எண்ணிக்கை. எனவே இப்போது செய்ய வேண்டியதெல்லாம்
~$\TMstroke$-களின் முதல் தொகுதியைத் தவிர மற்ற அனைத்தையும் அழித்துவிட்டு
நிற்பதே.

பின்வருவது ஒரு குறிப்பிடத்தக்க முடிவு:

\begin{thm}\ollabel{thm:universal-tm} உள்ளீடு $\tuple{e,n}$ இல்
  தொடங்கப்படும்போது பின்வருமாறு செயல்படும் ஓர்
  \emph{அனைத்தளாவிய டியூரிங் பொறி}~$U$ உள்ளது:
  \begin{enumerate}
    \item உள்ளீடு~$n$ இல் $M_e$ நின்றால், எனில் மட்டுமே அது நிற்கிறது, மேலும்
    \item $M_e$ வெளியீடு $m$ உடன் நின்றால்,~$U$ உம் அவ்வாறே நிற்கிறது.
  \end{enumerate}
  ஆகவே~$U$, பின்வருமாறு தரப்படும்
  $f\colon \Nat \times \Nat \pto \Nat$ என்ற சார்பைக் கணிக்கிறது:
  உள்ளீடு~$n$ இல் தொடங்கப்பட்ட $M_e$ வெளியீடு~$m$ உடன் நின்றால்
  $f(e,n) = m$; இல்லையெனில் அது வரையறுக்கப்படவில்லை.
\end{thm}

\begin{proof}
  ~அனைத்தளாவிய பொறி மிகவும் சிக்கலானது என்பதால்,~$U$ ஐ உண்மையில்
  உருவாக்குவது அடிப்படையில் இயலாதது. ஆனால் அது எவ்வாறு செயல்படுகிறது
  என்பதை வரைவாக விவரித்துவிட்டு, சர்ச்--டியூரிங் முன்வைப்பைப்
  பயன்படுத்தலாம். தொடக்கத்தில்~$U$ இன் நாடாவில் $e$
  எண்ணிக்கையிலான $\TMstroke$-களின் ஒரு தொகுதியைத் தொடர்ந்து $n$
  எண்ணிக்கையிலான~$\TMstroke$-களின் ஒரு தொகுதி உள்ளது. முதலில்
  சுட்டெண்~$e$ ஐ உள்ளீடு~$n$ கு வலப்புறத்தில் ``குறிவிலக்குகிறது''.
  இது பொறி~$M_e$ இன் கட்டளைகளை விவரிக்கும் எண்களின் ஒரு பட்டியலை
  (அதாவது,~$\TMblank$-களால் பிரிக்கப்பட்ட $\TMstroke$ தொகுதிகளை)
  உருவாக்குகிறது. அடுத்து~$U$,~$M_e$ இன் விவரிப்புக்கு வலப்புறத்தில்
  ~$M_e$ இன் தொடக்க நிலையின் எண்ணையும் எண்~$1$ ஐயும் நாடாவில் எழுதுகிறது.
  (மீண்டும், இவை ஒரும முறையில் $\TMstroke$ தொகுதிகளாகக்
  குறிப்பிடப்படுகின்றன.) அடுத்து, உள்ளீட்டை ($n$ எண்ணிக்கையிலான
  ~$\TMstroke$-களின் தொகுதியை) வலப்புறத்துக்கு நகலெடுக்கிறது---ஆனால்
  ஒவ்வொரு $\TMstroke$ ஐயும் மூன்று $\TMstroke$-களின் ஒரு தொகுதியால்
  மாற்றுகிறது ($\TMstroke$ குறியின் எண்~$3$ என்பதையும்,
  ~$\TMendtape$ இன் எண் $1$,~$\TMblank$ இன் எண் $2$ என்பதையும்
  நினைவுகூர்க). இத்தொகுதிகளின் தொடரின் இடது முனையில் (~$U$ இன்
  நாடாவில் $\TMblank$ குறிகளால் பிரிக்கப்பட்டது),~$\TMendtape$
  க்கான குறிமுறையான ஒரே ஒரு~$\TMstroke$ ஐ எழுதுகிறது.

  இப்போது~$U$ இன் நாடாவில் பின்வருவன உள்ளன: சுட்டெண்~$e$, எண்~$n$,
  தொடக்க நிலையின் குறியெண் (``நடப்பு நிலை''), தொடக்கத் தலை இருப்பிடத்தின்
  எண்~$1$ (``நடப்புத் தலை இருப்பிடம்''), ``நாடாவின்'' தொடக்க
  உள்ளடக்கம் (~$M_e$ இன் குறிகளின் குறியெண்களைக் குறிக்கும்
  ~$\TMstroke$ தொகுதிகளின் ஒரு தொடர்---அவை ``குறிகள்''---மேலும் அவை
  ~$\TMblank$-களால் பிரிக்கப்பட்டுள்ளன).

  உள்ளீடு~$n$ இல் தொடங்கப்பட்டால் $M_e$ என்ன செய்யுமோ அதை இப்போது
  ~பின்வருவனவற்றைச் செய்வதன் மூலம் அது உருவகப்படுத்துகிறது:
  \begin{enumerate}
    \item ``நடப்புத் தலை இருப்பிடத்தின்'' எண்~$k$ ஐக் கண்டறிக
    (தொடக்கத்தில் அது~$1$),
    \item அங்குள்ள ``குறி'' என்ன என்பதைக் காண ``நாடாவின்''
    $k$-ஆம் தொகுதிக்குச் செல்க,
    \item\ollabel{find-inst}%
    நடப்பு ``நிலைக்கும்'' ``குறிக்கும்'' பொருந்தும் கட்டளையைக் கண்டறிக,
    \item ``நாடாவின்'' $k$-ஆம் தொகுதிக்குத் திரும்பிச் சென்று, அங்குள்ள
    ``குறியை''~$M_e$ எழுதும் குறியின் குறியெண்ணால் மாற்றுக,
    \item நடப்பு ``நிலையைப்'' பதிவுசெய்யும் இடத்துக்குத் தலையை நகர்த்தி,
    அங்குள்ள எண்ணைப் புதிய நிலையின் எண்ணால் மாற்றுக,
    \item ``நாடா இருப்பிடத்தைப்'' பதிவுசெய்யும் இடத்துக்குச் சென்று, ஒரு
    ~$\TMstroke$ ஐ அழிக்க அல்லது ஒரு~$\TMstroke$ ஐச் சேர்க்க
    (கட்டளை முறையே இடப்புறம் அல்லது வலப்புறம் நகருமாறு கூறினால்).
    \item மீண்டும் செய்க.\footnote{இங்குள்ள சில நுணுக்கமான
    இடர்ப்பாடுகளை நாம் மேலோட்டமாகக் கடக்கிறோம். எடுத்துக்காட்டாக,
    ``நடப்புத் தலை இருப்பிடத்தை'' அது கண்காணிக்கும் எண்ணியைப் பெருக்கும்போது
    ~$U$ குக் கூடுதல் இடம் தேவைப்படலாம்---அந்நேர்வில் முழு ``நாடாவையும்''
    அது வலப்புறம் நகர்த்த வேண்டும்.}
  \end{enumerate}
  உள்ளீடு~$n$ இல் தொடங்கப்பட்ட $M_e$ ஒருபோதும் நிற்காவிட்டால்,
  ~$U$ உம் ஒருபோதும் நிற்பதில்லை; எனவே அதன் வெளியீடு
  வரையறுக்கப்படவில்லை.

  படி~\olref{find-inst} இல், நடப்பு ``நிலை''/``குறி'' ஜோடிக்கான
  எந்தக் கட்டளையும்~$M_e$ இன் விவரிப்பில் இல்லை என்று தெரியவந்தால்,
  ~$M_e$ நிற்கும். இது நிகழ்ந்தால்,~$U$ தனது நாடாவில் ``நாடாவுக்கு''
  இடப்புறமுள்ள பகுதியை அழிக்கிறது. மூன்று~$\TMstroke$-களின் ஒவ்வொரு
  தொகுதிக்கும் (~$M_e$ இன் நாடாவில் ஒரு~$\TMstroke$ ஐக் குறிக்கிறது),
  அது தனது சொந்த நாடாவின் இடது முனையில் ஒரு $\TMstroke$ ஐ எழுதிவிட்டு,
  ``நாடாவை'' ஒன்றன்பின் ஒன்றாக அழிக்கிறது. இது முடிந்ததும்,~$U$ இன்
  நாடாவில் நீளம்~$m$ ஆக உள்ள $\TMstroke$-களின் ஒரே தொகுதி உள்ளது.
  
  ``நாடாவில்'' மூன்று~$\TMstroke$-களின் ஒரு தொகுதி அல்லாத வேறொன்றை
  ~$U$ எதிர்கொண்டால், உடனடியாக நிற்கிறது. இந்நேர்வில்~$U$ இன்
  நாடாவில்~$\TMstroke$-களின் ஒரே தொகுதி இல்லாததால், அதன் வெளியீடு
  ஓர் இயல் எண் அல்ல; அதாவது இந்நேர்வில் $f(e,n)$
  வரையறுக்கப்படவில்லை.
\end{proof}

\end{document}
